\cleardoublepage
\phantomsection
\begin{planbox}[แผนบริหารการสอนปฏิบัติการที่ 1\quad การวัดละเอียดและการสอบเทียบไมโครปิเปต (สัปดาห์ที่ 1)]
\noindent\textbf{1. เวลาเรียนปฏิบัติการ}\quad 3 ชั่วโมง (180 นาที)\quad \textbf{ผลลัพธ์การเรียนรู้ที่ประเมิน}\quad CLO1, CLO3, CLO4

\noindent\textbf{2. วัตถุประสงค์เชิงพฤติกรรมประจำปฏิบัติการ}
\begin{enumerate}
    \item อ่านค่าและบันทึกผลการวัดมิติของจานเพาะเชื้อและหลอดทดลองด้วยเวอร์เนียร์คาลิปเปอร์และไมโครมิเตอร์พร้อมแปลงหน่วยเป็นไมโครเมตร ($\mu\text{m}$) ได้ถูกต้อง \textbf{(CLO1 เชื่อมโยง PLO2)}
    \item ปฏิบัติการดูดปล่อยของเหลวด้วยไมโครปิเปต (Micropipette) โดยใช้จังหวะการกดปุ่มที่ถูกต้อง และตรวจสอบความแม่นยำด้วยวิธีชั่งมวลน้ำกลั่นบนเครื่องชั่งดิจิทัลได้ \textbf{(CLO2, CLO3 เชื่อมโยง PLO1, PLO3)}
    \item บันทึกข้อมูลการทดลองตามความเป็นจริง คำนวณค่าเฉลี่ยและเปอร์เซ็นต์ความคลาดเคลื่อน และทำงานร่วมกับสมาชิกในกลุ่มได้อย่างมีประสิทธิภาพ \textbf{(CLO3, CLO4 เชื่อมโยง PLO1, PLO5, PLO6)}
\end{enumerate}

\noindent\textbf{3. แผนการจัดเวลาในคาบปฏิบัติการ 3 ชั่วโมง}
\begin{itemize}
    \item \textbf{ชั่วโมงที่ 1 (นาทีที่ 1--60)}\quad ชี้แจงกฎความปลอดภัย สาธิตการอ่านสเกลเวอร์เนียร์ ไมโครมิเตอร์ และการใช้ปุ่มกด 2 จังหวะของไมโครปิเปต
    \item \textbf{ชั่วโมงที่ 2 (นาทีที่ 61--120)}\quad นักศึกษาปฏิบัติการตอนที่ 1 (วัดมิติเครื่องแก้วจุลชีววิทยา) และตอนที่ 2 (สอบเทียบความแม่นยำไมโครปิเปต P200 และ P1000)
    \item \textbf{ชั่วโมงที่ 3 (นาทีที่ 121--180)}\quad คำนวณค่าเฉลี่ยและเปอร์เซ็นต์ความคลาดเคลื่อน อภิปรายผลร่วมกัน และตรวจเช็กความสะอาดของอุปกรณ์ก่อนส่งใบบันทึกผล
\end{itemize}
\end{planbox}
\clearpage

\chapter{การวัดละเอียดและการสอบเทียบไมโครปิเปต}
\index{เวอร์เนียร์คาลิปเปอร์}\index{ไมโครมิเตอร์}\index{การสอบเทียบไมโครปิเปต}\index{เครื่องชั่งดิจิทัล}

\section{หลักการทางฟิสิกส์และทฤษฎีเบื้องหลังการทดลอง}

ในงานปฏิบัติการจุลชีววิทยา ความแม่นยำของการวัดปริมาตรสารละลายระดับไมโครลิตร ($\mu\text{L}$) มีผลโดยตรงต่อความถูกต้องของการเจือจางเชื้อแบบลำดับส่วน (\textbf{Serial Dilution}) และการเตรียมปฏิกิริยาระดับโมเลกุล เครื่องมือหลักที่ใช้คือ \textbf{ไมโครปิเปต} (\textbf{Micropipette}) ซึ่งอาศัยหลักการแทนที่อากาศ (\textbf{Air Displacement}) ของลูกสูบภายในกระบอกสูบ

วิธีมาตรฐานสากลในการตรวจสอบว่าไมโครปิเปตยังจ่ายปริมาตรได้แม่นยำหรือไม่ คือ \textbf{วิธีชั่งมวลน้ำกลั่น} (\textbf{Gravimetric Calibration}) โดยอาศัยหลักความหนาแน่นของน้ำบริสุทธิ์ที่อุณหภูมิห้อง ($\rho = m / V \approx 1.000\text{ g/mL}$ หรือ $0.001\text{ g/}\mu\text{L}$) กล่าวคือ \textbf{น้ำกลั่นปริมาตร $1\text{ }\mu\text{L}$ จะมีมวลเท่ากับ $0.0010\text{ g}$ ($1\text{ mg}$) พอดี} ดังนั้น น้ำกลั่นปริมาตร $1{,}000\text{ }\mu\text{L}$ จึงมีมวลมาตรฐานเท่ากับ $1.0000\text{ g}$ และน้ำกลั่นปริมาตร $200\text{ }\mu\text{L}$ จึงมีมวลมาตรฐานเท่ากับ $0.2000\text{ g}$

\begin{conceptbox}[สูตรการแปลงหน่วยและการหาเปอร์เซ็นต์ความคลาดเคลื่อน]
\begin{enumerate}
    \item \textbf{การแปลงหน่วยความยาวและปริมาตร}\quad $1\text{ mm} = 1{,}000\text{ }\mu\text{m}$ และ $1\text{ mL} = 1{,}000\text{ }\mu\text{L}$
    \item \textbf{การอ่านค่าเวอร์เนียร์คาลิปเปอร์}\quad ค่าที่อ่านได้ $=$ สเกลหลัก ($\text{mm}$) $+$ (ขีดสเกลเลื่อนที่ตรงกัน $\times$ ความละเอียดเครื่องมือ)
    \item \textbf{เปอร์เซ็นต์ความคลาดเคลื่อนของไมโครปิเปต}
    \begin{equation}
        \text{ความคลาดเคลื่อน (\%)} = \frac{|\bar{m}_{\text{ชั่งได้จริง}} - m_{\text{มาตรฐาน}}|}{m_{\text{มาตรฐาน}}} \times 100\%
    \end{equation}
\end{enumerate}
\end{conceptbox}

\begin{worked}[การคำนวณความคลาดเคลื่อนของการดูดน้ำกลั่นด้วยไมโครปิเปต P200]
\textbf{โจทย์ตัวอย่าง}\quad นักศึกษาตั้งค่าไมโครปิเปต P200 ที่ปริมาตร $200\text{ }\mu\text{L}$ (มวลน้ำมาตรฐานเท่ากับ $0.2000\text{ g}$) เมื่อดูดน้ำกลั่นหยดลงบนเครื่องชั่งดิจิทัล 3 ครั้ง ได้ค่ามวลเท่ากับ $0.1960\text{ g}$, $0.1980\text{ g}$ และ $0.1970\text{ g}$ จงคำนวณมวลเฉลี่ยและเปอร์เซ็นต์ความคลาดเคลื่อน

\textbf{วิธีคำนวณทีละขั้นตอน}
\begin{enumerate}
    \item \textbf{หาค่ามวลเฉลี่ย}\quad $\bar{m} = \dfrac{0.1960 + 0.1980 + 0.1970}{3} = 0.1970\text{ g}$ (คิดเป็นปริมาตรจริง $197.0\text{ }\mu\text{L}$)
    \item \textbf{หาเปอร์เซ็นต์ความคลาดเคลื่อน}\quad $\dfrac{|0.1970 - 0.2000|}{0.2000} \times 100\% = \dfrac{0.0030}{0.2000} \times 100\% = 1.50\%$
\end{enumerate}
\end{worked}

\section{แผนภาพชุดอุปกรณ์การทดลองและข้อควรระวังความปลอดภัย}

\begin{figure}[H]
\centering
\resizebox{0.86\textwidth}{!}{%
\begin{tikzpicture}[font=\sffamily, >=Stealth]
    % Left: Micropipette 2-Stop Mechanism
    \begin{scope}[shift={(0,0)}]
        \draw[fill=boxBlueBg, draw=physRoyalBlue, line width=1.2pt, rounded corners=4pt] (0,0) rectangle (6.8,4.4);
        \node[font=\bfseries\color{physRoyalBlue}] at (3.4,3.95) {(ก) จังหวะการกดปุ่มไมโครปิเปต};
        % Plunger positions
        \fill[physCyan!30, draw=physRoyalBlue, line width=1.1pt] (1.0,1.2) rectangle (2.2,2.8);
        \fill[physAmber, draw=charcoal, line width=1.0pt] (1.3,2.8) rectangle (1.9,3.4);
        \node[below, font=\scriptsize\bfseries] at (1.6,1.15) {ปล่อยปกติ};

        \fill[physCyan!30, draw=physRoyalBlue, line width=1.1pt] (2.8,1.2) rectangle (4.0,2.8);
        \fill[bioEmerald, draw=charcoal, line width=1.0pt] (3.1,2.8) rectangle (3.7,3.1);
        \node[below, font=\scriptsize\bfseries, color=bioEmerald!90!black, align=center] at (3.4,1.15) {จังหวะที่ 1\\(ดูดสาร)};

        \fill[physCyan!30, draw=physRoyalBlue, line width=1.1pt] (4.6,1.2) rectangle (5.8,2.8);
        \fill[crimsonAlert, draw=charcoal, line width=1.0pt] (4.9,2.8) rectangle (5.5,2.92);
        \node[below, font=\scriptsize\bfseries, color=crimsonAlert, align=center] at (5.2,1.15) {จังหวะที่ 2\\(ไล่หยดสุดท้าย)};
    \end{scope}

    % Right: Digital Balance Gravimetric Setup
    \begin{scope}[shift={(7.6,0)}]
        \draw[fill=boxGreenBg, draw=bioEmerald, line width=1.2pt, rounded corners=4pt] (0,0) rectangle (6.8,4.4);
        \node[font=\bfseries\color{bioEmerald!90!black}] at (3.4,3.95) {(ข) การชั่งมวลน้ำกลั่นสอบเทียบปิเปต};
        % Balance Base
        \fill[charcoal!85, rounded corners=2pt] (1.2,0.6) rectangle (5.6,1.4);
        \node[fill=bioEmerald!30!white, text=charcoal, font=\small\bfseries, rounded corners=1pt, inner sep=2pt] at (3.4,1.0) {$0.5000\text{ g}$};
        \draw[physSlate, line width=1.4pt] (2.2,1.4) -- (4.6,1.4);
        % Weighing beaker with water
        \fill[physCyan!25] (2.6,1.45) rectangle (4.2,2.3);
        \draw[physRoyalBlue, line width=1.2pt] (2.6,2.7) -- (2.6,1.45) -- (4.2,1.45) -- (4.2,2.7);
        % Pipette tip dropping water
        \fill[physAmber!70, draw=physAmber!90!black, line width=1.0pt] (3.25,3.5) -- (3.55,3.5) -- (3.43,2.75) -- (3.37,2.75) -- cycle;
        \fill[physCyan] (3.4,2.55) circle (2.2pt);
    \end{scope}
\end{tikzpicture}%
}
\caption{กลไกปุ่มกด 2 จังหวะของไมโครปิเปตและการจัดชุดสอบเทียบปริมาตรด้วยเครื่องชั่งดิจิทัล}
\end{figure}

\begin{equipbox}[รายการอุปกรณ์สำหรับแต่ละกลุ่ม (กลุ่มละ 3--4 คน)]
\begin{itemize}
    \item เวอร์เนียร์คาลิปเปอร์ (ความละเอียด $0.05\text{ mm}$) จำนวน 1 อัน และไมโครมิเตอร์ (ความละเอียด $0.01\text{ mm}$) จำนวน 1 อัน
    \item จานเพาะเชื้อแก้วปลอดเชื้อ 1 คู่, หลอดทดลอง 1 หลอด และแผ่นสไลด์พร้อมแผ่นกระจกปิดสไลด์ (Coverslip)
    \item ไมโครปิเปตขนาด P200 ($20\text{--}200\text{ }\mu\text{L}$) และ P1000 ($100\text{--}1{,}000\text{ }\mu\text{L}$) พร้อมกล่องทิปพลาสติก
    \item เครื่องชั่งดิจิทัลทศนิยม 4 ตำแหน่ง บีกเกอร์พลาสติกขนาดเล็ก และขวดน้ำกลั่นที่วัดอุณหภูมิห้องแล้ว
\end{itemize}
\end{equipbox}

\begin{warningbox}[ข้อห้ามเด็ดขาดในการใช้งานไมโครปิเปต]
\begin{enumerate}
    \item \textbf{ห้ามหมุนปรับตัวเลขปริมาตรเกินช่วงที่ระบุไว้บนตัวปิเปต} (เช่น ห้ามหมุน P200 เกิน $200\text{ }\mu\text{L}$ หรือต่ำกว่า $20\text{ }\mu\text{L}$) เพราะเฟืองสปริงภายในจะพังเสียหายทันที
    \item \textbf{ห้ามวางไมโครปิเปตที่มีของเหลวค้างในปลายทิปนอนราบลงบนโต๊ะ} เพราะของเหลวจะไหลย้อนเข้าไปกัดกร่อนลูกสูบโลหะด้านใน
\end{enumerate}
\end{warningbox}

\section{ขั้นตอนการปฏิบัติการทดลอง}

\subsection{ตอนที่ 1 การวัดมิติของภาชนะเพาะเลี้ยงเชื้อจุลินทรีย์}
\begin{enumerate}
    \item ตรวจสอบค่าศูนย์ของเวอร์เนียร์คาลิปเปอร์และไมโครมิเตอร์ก่อนเริ่มวัด
    \item ใช้เขี้ยววัดด้านในของเวอร์เนียร์คาลิปเปอร์วัด\textbf{เส้นผ่านศูนย์กลางภายใน}ของจานเพาะเชื้อแก้ว และใช้ก้านวัดความลึกวัด\textbf{ความลึกภายใน}ของจานเพาะเชื้อ ทำซ้ำ 3 ตำแหน่ง บันทึกค่าในหน่วยมิลลิเมตร ($\text{mm}$) แล้วแปลงเป็นไมโครเมตร ($\mu\text{m}$)
    \item ใช้ไมโครมิเตอร์วัด\textbf{ความหนาของแผ่นสไลด์แก้ว} และ\textbf{ความหนาของแผ่นกระจกปิดสไลด์} (Coverslip) อย่างละ 3 ครั้ง โดยหมุนที่ปุ่มปรับแรงดันท้ายแกนจนได้ยินเสียงคลิกเบาๆ 3 ครั้ง แล้วอ่านค่า
\end{enumerate}

\subsection{ตอนที่ 2 การสอบเทียบความแม่นและความเที่ยงของไมโครปิเปต}
\begin{enumerate}
    \item วางบีกเกอร์พลาสติกใบเล็กบนจานเครื่องชั่งดิจิทัล ปิดกระจกกันลม แล้วกดปุ่มเซตศูนย์ (\textbf{Tare} ให้หน้าจอแสดง $0.0000\text{ g}$)
    \item ปรับตั้งค่าไมโครปิเปต P200 ไปที่ $100\text{ }\mu\text{L}$ สวมทิปให้แน่น กดปุ่มลงถึง\textbf{จังหวะที่ 1} จุ่มปลายทิปลงใต้ผิวน้ำกลั่นประมาณ $2\text{--}3\text{ mm}$ แล้วค่อยๆ ผ่อนนิ้วขึ้นช้าๆ เพื่อดูดน้ำกลั่นขึ้นมา
    \item นำปลายทิปแตะข้างผนังบีกเกอร์บนเครื่องชั่ง กดปุ่มลงถึงจังหวะที่ 1 แล้วกดต่อไปถึง\textbf{จังหวะที่ 2} เพื่อไล่หยดน้ำออกให้หมด ปิดกระจกเครื่องชั่ง อ่านค่ามวลเป็นกรัม ($\text{g}$) บันทึกผล กดปุ่ม Tare แล้วทำซ้ำให้ครบ 3 ครั้ง
    \item เปลี่ยนเป็นไมโครปิเปต P1000 ตั้งค่าที่ $500\text{ }\mu\text{L}$ และ $1{,}000\text{ }\mu\text{L}$ ทำการทดลองซ้ำเช่นเดียวกัน
\end{enumerate}

\section{ใบบันทึกผลการทดลองประจำกลุ่ม}

\begin{worksheetbox}[ใบบันทึกผลการทดลองที่ 1\quad การวัดละเอียดและการสอบเทียบไมโครปิเปต]
\noindent\textbf{กลุ่มปฏิบัติการที่} \dotfill \textbf{วันที่ทำการทดลอง} \dotfill \textbf{อุณหภูมิน้ำกลั่น} \dotfill $^\circ\text{C}$\\[4pt]
\noindent\textbf{รายชื่อสมาชิกกลุ่ม} 1. \dotfill 2. \dotfill 3. \dotfill

\vspace{6pt}
\noindent\textbf{ตารางบันทึกผลตอนที่ 1 การวัดมิติภาชนะและแผ่นสไลด์ด้วยเวอร์เนียร์และไมโครมิเตอร์}
\begin{center}
\footnotesize\setlength{\tabcolsep}{3.5pt}
\begin{tabularx}{\linewidth}{>{\raggedright\arraybackslash}X c c c c c}
\toprule
\textbf{รายการที่วัด} & \textbf{ครั้งที่ 1} & \textbf{ครั้งที่ 2} & \textbf{ครั้งที่ 3} & \textbf{ค่าเฉลี่ย ($\text{mm}$)} & \textbf{แปลงเป็นหน่วย $\mu\text{m}$} \\
\midrule
เส้นผ่านศูนย์กลางในจานเพาะเชื้อ ($\text{mm}$) & ............ & ............ & ............ & ............ & \dotfill \\
ความลึกภายในจานเพาะเชื้อ ($\text{mm}$) & ............ & ............ & ............ & ............ & \dotfill \\
ความหนาแผ่นสไลด์แก้ว ($\text{mm}$) & ............ & ............ & ............ & ............ & \dotfill \\
ความหนากระจกปิดสไลด์ ($\text{mm}$) & ............ & ............ & ............ & ............ & \dotfill \\
\bottomrule
\end{tabularx}
\end{center}

\vspace{4pt}
\noindent\textbf{ตารางบันทึกผลตอนที่ 2 การสอบเทียบความแม่นยำของไมโครปิเปตด้วยวิธีชั่งมวลน้ำกลั่น}
\begin{center}
\footnotesize\setlength{\tabcolsep}{3.5pt}
\begin{tabularx}{\linewidth}{>{\raggedright\arraybackslash}X c c c c c}
\toprule
\textbf{ขนาดปิเปตและปริมาตรที่ตั้ง} & \textbf{ครั้งที่ 1 ($\text{g}$)} & \textbf{ครั้งที่ 2 ($\text{g}$)} & \textbf{ครั้งที่ 3 ($\text{g}$)} & \textbf{เฉลี่ย ($\text{g}$)} & \textbf{คลาดเคลื่อน (\%)} \\
\midrule
P200 ตั้งที่ $100\text{ }\mu\text{L}$ ($0.1000\text{ g}$) & ........ & ........ & ........ & ........ & \dotfill \\
P1000 ตั้งที่ $500\text{ }\mu\text{L}$ ($0.5000\text{ g}$) & ........ & ........ & ........ & ........ & \dotfill \\
P1000 ตั้งที่ $1{,}000\text{ }\mu\text{L}$ ($1.0000\text{ g}$) & ........ & ........ & ........ & ........ & \dotfill \\
\bottomrule
\end{tabularx}
\end{center}

\vspace{4pt}
\noindent\textbf{แสดงวิธีคำนวณตัวอย่าง (เลือกจากตารางตอนที่ 2 จำนวน 1 รายการ)}\\[4pt]
\dotfill\\[6pt]
\dotfill

\noindent\textbf{สรุปและอภิปรายผลการทดลอง (วิเคราะห์สาเหตุความคลาดเคลื่อนที่พบจริง)}\\[4pt]
\dotfill\\[6pt]
\dotfill
\end{worksheetbox}

\section*{คำถามท้ายปฏิบัติการที่ 1}
\begin{enumerate}
    \item หากนักศึกษากดปุ่มไมโครปิเปตลงไปจนสุดถึงจังหวะที่ 2 ก่อนจุ่มปลายทิปดูดสารละลายเชื้อแบคทีเรีย มวลของเหลวที่ชั่งได้ในตอนที่ 2 จะมากกว่าหรือน้อยกว่าค่ามาตรฐาน เพราะเหตุใด
    \item แผ่นกระจกปิดสไลด์ (Coverslip) ที่วัดความหนาได้ในตอนที่ 1 มีความหนากี่ไมโครเมตร ($\mu\text{m}$) และมีความสำคัญอย่างไรต่อการป้องกันหน้าเลนส์ใกล้วัตถุของกล้องจุลทรรศน์
\end{enumerate}
